\section{Basic Category Theory}\label{secA}

The main prerequisite for reading this paper is knowledge of basic category theory. As such, we include this appendix to serve as a reference for some important concepts which are of importance to the paper. We indent this appendix purely as a \emph{reference} for those who already know category theory, and simply want a refresher on some concepts. In particular, we do not pretend that it is possible to learn category theory from this appendix, and although picking up one or two new ideas may be possible, we prefer that the reader first familiarize themselves with basic category theory using another text.

To this end, we recommend \cite{ctx} and \cite{lein} for the advanced undergraduate or early graduate. We also note that \cite{alg0} provides a slower and more motivated approach to category theory, since it builds algebra using category theoretic reasoning. However, this is mostly an algebra book, and does not go far into the weeds of category theory; we recommend it to less advanced students as a precursor to a full text on the topic, since it provides material that is crucial to understand many of the examples in other texts, and helps motivate the earlier sections of these texts. We note \cite{ml} is a solid reference on the topic, but we discourage learning from it, and we also dislike Mac Lane's notational conventions.

\subsection{Categories, functors, natural transformations}\label{secA.1}
\begin{definition}\label{A.1.1}
    A \emph{category} $\C$ is a class $\Obj(\C)$ of \emph{objects} of $\C$, together with a class $\Mor(\C)$ of \emph{morphisms} of $\C$, and four functions:
    \begin{itemize}
        \item A function $\dom : \Mor(\C) \to \Obj(\C)$, assigning each morphism $f\in\Mor(\C)$ a \emph{domain} $\dom(f)$;
        \item A function $\cod : \Mor(\C)\to\Obj(\C)$, assigning each morphism $f\in\Mor(\C)$ a \emph{codomain} $\cod(f)$;
        \item A function $1_\bullet : \Obj(\C)\to\Mor(\C)$, assigning each object $x\in\Obj(\C)$ a \emph{identity morphism} $1_x$;
        \item If $f,g\in\Mor(\C)$ are morphisms such that $\cod(f)=\dom(g)$, we say that the pair $(g,f)$ is a \emph{composable pair}. There is then a function $\circ$ called \emph{composition} from the set of composable pairs to $\Mor(\C)$
        \[
        \circ : \{(g,f)\in\Mor(\C)\times\Mor(\C) \mid \dom(g)=\cod(f)\} \to \Mor(\C)
        \]
        that takes each composable pair $(g,f)$, and maps it to a morphism $g\circ f$, called the \emph{composite} of $g$ and $f$.
    \end{itemize}
    These functions must satisfy the following axioms:
    \begin{enumerate}
        \item For any $x\in\Obj(\C)$, the identity morphism $1_x$ must have $\dom(1_x)=\cod(1_x)=x$.
        \item (right unitor) For any $f\in\Mor(\C)$ with $\dom(f)=x$, we must have the strict equality $f\circ 1_x = f$.
        \item (left unitor) For any $g\in\Mor(\C)$ with $\cod(g)=x$, we must have the strict equality $1_x\circ g= g$.
        \item (associativity) For any two composable pairs $(h,g)$ and $(g,f)$, we must have an equality $(h\circ g)\circ f=h\circ(g\circ f)$. We call $(h,g,f)$ a \emph{composable triple}, and write the composite $h\circ g\circ f$.
    \end{enumerate}
\end{definition}

\begin{notation}\label{A.1.2}
    If $f\in\Mor(\C)$ with $\dom(f)=x$ and $\cod(f)=y$, we say that $f$ \emph{is a morphism from $x$ to $y$}, and write $f : x \to y$ or $x\xlongrightarrow{f} y$. We write $\C(x,y)$ to denote the class of all morphisms from $x$ to $y$. If $(g,f)$ is a composable pair, we often denote the composite $g\circ f$ by juxtaposition $gf$, omitting the $\circ$ where possible. We often drop the $\Obj$, and write $x\in\C$ to denote that $x$ is an object of $\C$. We typically will never write $f\in\Mor(\C)$, and prefer to say \emph{$f$ is a morphism of $\C$} using language.
\end{notation}

\begin{remark}\label{A.1.3}
    Most categories in nature are \emph{locally small}, meaning for any objects $x,y\in\C$, the class $\C(x,y)$ is in fact a \emph{set}. We will only need to consider locally small categories in this paper, so we may safely consider $\C(x,y)$ to be a set throughout. Some authors use the notation $\Hom(x,y)$ or $\Hom_\C(x,y)$ for the set $\C(x,y)$, so we often refer to $\C(x,y)$ as a \emph{hom-set}. Most of what we prove here will work in a more general setting where $\C(x,y)$ is a proper class, but working through the foundational issues adds unnecessary complexity, and some authors even go as far as to require that hom-sets indeed be sets in their definition of a category. When the class of objects of a category $\C$ is a set, we call the category \emph{small}. If a category is not small, we call it \emph{large}.
\end{remark}

\begin{example}\label{A.1.4}
    A partial order on a set $P$ is a relation $\preceq$ which is
    \begin{enumerate}
        \item reflexive: for all $x\in P$, $x\preceq x$;
        \item transitive: for all $x,y,z\in P$, if $x\preceq y$ and $y\preceq z$, then $x\preceq z$;
        \item antisymmetric: for all $x,y\in P$, $x\preceq y$ and $y\preceq x$ \emph{if and only if} $x=y$.
    \end{enumerate}
    We can encode the information this relation carries into a category. If $x\preceq y$, then there is a unique morphism $x\to y$, and otherwise there is no morphism. Such a category is called an \emph{poset category}. A common type of poset category is for ordinal numbers: any ordinal number $n$, when viewed as a set, forms a poset under the standard $\leq$ relation. We thus can regard $n$ as a poset category under this relation, and typically write $\mathbf{n}$ for this category.
\end{example}

\begin{definition}\label{A.1.5}
    Let $\C$ and $\D$ be categories. A \emph{functor} $F : \C \to \D$ consists of
    \begin{itemize}
        \item an object function $\Obj(\C)\to\Obj(\D)$,
        \item and a morphism function $\Mor(\C)\to\Mor(\D)$,
    \end{itemize}
    both of which are typically denoted by $F$. These functions must preserve domain, codomain, identity, and composition:
    \begin{enumerate}
        \item if $f : x\to y$ in $\C$, then $F(f) : F(x)\to F(y)$ in $\D$;
        \item for any composable pair $(g,f)$ in $\C$, there is an equality $F(g\circ f)=F(g)\circ F(f)$;
        \item for any object $x\in\C$, $F(1_x)=1_{F(x)}$.
    \end{enumerate}
\end{definition}

\begin{remark}\label{A.1.6}
    It's shown fairly easily that given a commutative diagram, if we apply a functor $F$ to every object and every morphism in the diagram, it will still commute. For example, if the square
    \[\begin{tikzcd}
        a\ar["f", r]\ar["g"', d] & b \ar["h", d]\\
        c \ar["k"', r] & d
    \end{tikzcd}\]
    commutes, then so does the square
    \[\begin{tikzcd}
        F(a)\ar["F(f)", r]\ar["F(g)"', d] & F(b) \ar["F(h)", d]\\
        F(c) \ar["F(k)"', r] & F(d).
    \end{tikzcd}\]
\end{remark}

\begin{terminology}\label{A.1.7}
    We will usually refer to the properties (1), (2), and (3) of \cref{A.1.5}, and more generally \cref{A.1.6}, as \emph{functoriality}.
\end{terminology}

\begin{definition}\label{A.1.8}
    Let $\C,\D$ be categories, and $F,G : \C\to\D$ be functors. A \emph{natural transformation} $\alpha : F \Rightarrow G$ is a function which assigns each $x\in\C$ to a morphism $\alpha_x : F(x)\to G(x)$ in $\D$, such that for any $f : x\to y$ in $\C$, the diagram
    \[\begin{tikzcd}
        F(x) \ar["F(f)", r]\ar["\alpha_x"', d] & F(y)\ar["\alpha_y", d]\\
        G(x)\ar["G(f)"', r] & G(y)
    \end{tikzcd}\]
    commutes.
\end{definition}

\begin{remark}\label{A.1.9}
    We obviously note that we can compose funtors in the same way we compose functions. We can also compose natural transformations: if $F,G,H : \C\to\D$ are functors and $\alpha : F\Rightarrow G$, $\beta : G\Rightarrow H$ are natural transformations, then the \emph{vertical composite} of $\beta$ and $\alpha$ is a natural transormation $\beta\circ\alpha : F \Rightarrow H$, with components given by $(\beta\circ\alpha)_x=\beta_{x}\circ\alpha_x$.
\end{remark}

\begin{notation}\label{A.1.10}
We have various conventions for diagrams including natural transformations. A simple diagram
\[\begin{tikzcd}
	\C & \D
	\arrow[""{name=0, anchor=center, inner sep=0}, "F", curve={height=-12pt}, from=1-1, to=1-2]
	\arrow[""{name=1, anchor=center, inner sep=0}, "G"', curve={height=12pt}, from=1-1, to=1-2]
	\arrow["\alpha", shorten=2pt, Rightarrow, from=0, to=1]
\end{tikzcd}\]
depicts that $F$ and $G$ are functors, and there is a natural transformation $\alpha$ between them. More complex diagrams such as
\[\begin{tikzcd}
	\C && \mathcal{E} \\
	& \D
	\arrow[""{name=0, anchor=center, inner sep=0}, "F", from=1-1, to=1-3]
	\arrow["G"', from=1-1, to=2-2]
	\arrow["H"', from=2-2, to=1-3]
	\arrow[shorten=2pt, Rightarrow, from=0, to=2-2, "\alpha"]
\end{tikzcd}\]
denote simply that $\alpha$ is a natural transformation from $F$ to the composite $H\circ G$.
\end{notation}

\begin{construction}\label{A.1.11}
    \Cref{A.1.9} provides us a way to compose natural transformations along functors. However, we can also compose natural transformations along categories, using what is called \emph{horizontal composition}. Given a diagram
    \[\begin{tikzcd}
	\C & \D & \E,
	\arrow[""{name=0, anchor=center, inner sep=0}, "{F_1}", curve={height=-12pt}, from=1-1, to=1-2]
	\arrow[""{name=1, anchor=center, inner sep=0}, "{G_1}"', curve={height=12pt}, from=1-1, to=1-2]
	\arrow[""{name=2, anchor=center, inner sep=0}, "{F_2}", curve={height=-12pt}, from=1-2, to=1-3]
	\arrow[""{name=3, anchor=center, inner sep=0}, "{G_2}"', curve={height=12pt}, from=1-2, to=1-3]
	\arrow["\alpha", shorten=2pt, Rightarrow, from=0, to=1]
	\arrow["\beta", shorten=2pt, Rightarrow, from=2, to=3]
\end{tikzcd}\]
    the \emph{horizontal composite} of $\beta$ and $\alpha$, denoted $\beta\cdot\alpha$, is a natural transformation
    \[\begin{tikzcd}
	\C && \E
	\arrow[""{name=0, anchor=center, inner sep=0}, "{F_2\circ F_1}", curve={height=-12pt}, from=1-1, to=1-3]
	\arrow[""{name=1, anchor=center, inner sep=0}, "{G_2\circ G_1}"', curve={height=12pt}, from=1-1, to=1-3]
	\arrow["{\beta\cdot\alpha}", shorten=2pt, Rightarrow, from=0, to=1]
\end{tikzcd}\]
    The components $(\beta\cdot\alpha)_x$ can be defined as either of the following expressions:
    \[
    (\beta\cdot\alpha)_x = \beta_{G_1(x)}\circ F_2(\alpha_x) \qquad \qquad (\beta\cdot\alpha)_x = G_2(\alpha_x)\circ \beta_{F_1(x)}.
    \]
    These two composites are shown to be equal by \cref{A.1.13} and \cref{A.1.14}.
\end{construction}

\begin{notation}\label{A.1.12}
    As in \cref{A.1.2}, we often denote the composite $G\circ F$ of two functors by $GF$, omitting the $\circ$. Since there are two composites of natural transformations, notation becomes slightly more tricky. We prefer to never use the $\cdot$ notation for horizontal composition, so we will almost always denote the horizontal composite of $\beta$ and $\alpha$ by juxtaposition $\beta\alpha$. If both horizontal and vertical composition are being used, we will denote vertical composition by $\circ$. If only vertical composition is being used, we will usually denote it by juxtaposition. If only one type of composition is being used, we will usually clarify whether it is the horizontal or vertical composite, unless it is clear from context.
\end{notation}

\begin{remark}\label{A.1.13}
    Given a diagram
    \[\begin{tikzcd}
	\C & \D & \E,
	\arrow[""{name=0, anchor=center, inner sep=0}, "F", curve={height=-12pt}, from=1-1, to=1-2]
	\arrow[""{name=1, anchor=center, inner sep=0}, "G"', curve={height=12pt}, from=1-1, to=1-2]
	\arrow[""{name=2, anchor=center, inner sep=0}, "H", curve={height=-12pt}, from=1-2, to=1-3]
	\arrow[""{name=3, anchor=center, inner sep=0}, "H"', curve={height=12pt}, from=1-2, to=1-3]
	\arrow["\alpha", shorten=2pt, Rightarrow, from=0, to=1]
	\arrow[shorten=2pt, equals, from=2, to=3]
\end{tikzcd}\]
    we define the horizontal composite of $\alpha$ and the identity natural transformation $1_H$. Using either definition in \cref{A.1.11}, we obtain the definitions
    \[
    (1_H\alpha)_x = (1_H)_{G(x)}\circ H(\alpha_x)=H(\alpha_x)\circ (1_H)_{F(x)}.
    \]
    We have not yet shown the second equality holds, but we can show this fairly easily for this special case. Note that each component $(1_H)_x$ is just the identity $1_{H(x)}$. We can write these composites as the following diagram, which commutes if and only if the composites are indeed equal:
    \[\begin{tikzcd}
	{H(F(x))} & {H(G(x))} \\
	{H(F(x))} & {H(G(x))}
	\arrow["{H(\alpha_x)}", from=1-1, to=1-2]
	\arrow["{1_{H(F(x))}}"', from=1-1, to=2-1]
	\arrow["{1_{H(G(x))}}", from=1-2, to=2-2]
	\arrow["{H(\alpha_x)}"', from=2-1, to=2-2]
    \end{tikzcd}\]
    Note that this diagram is simply $H$ applied to the naturality condition for $\alpha$ with respect to the identity arrow, so this diagram commutes by \cref{A.1.8} and \cref{A.1.6}. Thus, the composites are indeed equal. The notation $(1_H)_x$ is very cumbersome, but luckily we usually only see identity natural transformations when they are being horizontally composed with something else. We thus adopt a new notation for horizontal composites with the identity morphism: $1_H\cdot \alpha$ is instead denoted as $H\alpha$. This notation is very nice, since $(H\alpha)_x=H(\alpha_x)$.

    We also see that the notation $\alpha H$ is nice. Consider the diagram
    \[\begin{tikzcd}
	\E & \C & \D,
	\arrow[""{name=0, anchor=center, inner sep=0}, "H", curve={height=-12pt}, from=1-1, to=1-2]
	\arrow[""{name=1, anchor=center, inner sep=0}, "H"', curve={height=12pt}, from=1-1, to=1-2]
	\arrow[""{name=2, anchor=center, inner sep=0}, "F", curve={height=-12pt}, from=1-2, to=1-3]
	\arrow[""{name=3, anchor=center, inner sep=0}, "G"', curve={height=12pt}, from=1-2, to=1-3]
	\arrow["\alpha", shorten=2pt, Rightarrow, from=2, to=3]
	\arrow[shorten=2pt, equals, from=0, to=1]
    \end{tikzcd}\]
    By the first formula in \cref{A.1.11}, the component at $x$ is given as
    \[
    (\alpha H)_x = \alpha_{H(x)}\circ F(H_x) = \alpha_{H(x)}\circ F(1_{H(x)})=\alpha_{H(x)}\circ 1_{FH(x)}=\alpha_{H(x)},
    \]
    where the subsequent steps follow by our notation $H_x=(1_H)_x=1_{H(x)}$, by \cref{A.1.5}, and by \cref{A.1.1}. By the second formula, we have
    \[
    (\alpha H)_x = G(H_x)\circ \alpha_{H(x)}=G(1_{H(x)})\circ\alpha_{H(x)}=1_{GH(x)}\circ\alpha_{H(x)}=\alpha_{H(x)},
    \]
    where the subsequent steps follow for the same reasons. We thus see that $(\alpha H)_x=\alpha_{H(x)}$, both justifying the equality in \cref{A.1.11} and justifying our notation.
\end{remark}

\begin{remark}\label{A.1.14}
    We can now justify the equality in \cref{A.1.11}. Recall the situation depicted by the diagram therein, and note that we can alternatively view the two formulas given as
    \[
    (\beta\alpha)_x=(\beta G_1)_x\circ (F_2\alpha)_x,\qquad (\beta\alpha)_x = (G_2\alpha)_x\circ (\beta F_1)_x.
    \]
    We thus have
    \[
    \beta\alpha = \beta G_1\circ F_2\alpha,\qquad \beta\alpha = G_2\alpha\circ \beta F_1.
    \]
    To show this indeed holds, let $x\in\C$, and note that $\alpha_x : F_1(x)\to G_1(x)$. By functoriality of $F_2$ and $G_2$, we thus have two morphisms
    \[\begin{tikzcd}[row sep=0em]
	{F_2F_1(x)} & {F_2G_1(x)} \\
	{G_2F_1(x)} & {G_2G_1(x).}
	\arrow["{F_2(\alpha_x)}", from=1-1, to=1-2]
	\arrow["{G_2(\alpha_x)}", from=2-1, to=2-2]
    \end{tikzcd}\]
    Since $\beta$ is natural from $F_2$ to $G_2$, by \cref{A.1.8} the square
    \[\begin{tikzcd}
	{F_2F_1(x)} & {F_2G_1(x)} \\
	{G_2F_1(x)} & {G_2G_1(x)}
	\arrow["{F_2(\alpha_x)}", from=1-1, to=1-2]
	\arrow["{\beta_{F_1(x)}}"', from=1-1, to=2-1]
	\arrow["{\beta_{G_1(x)}}", from=1-2, to=2-2]
	\arrow["{G_2(\alpha_x)}", from=2-1, to=2-2]
    \end{tikzcd}\]
    commutes. Recognize that this expresses the equality
    \[
    \beta_{G_1(x)}\circ F_2(\alpha_x)=G_2(\alpha_x)\circ \beta_{F_1(x)},
    \]
    which is, by the notation introduced in \cref{A.1.13}, the same as
    \[
    (\beta G_1)_x\circ (F_2\alpha)_x = (G_2\alpha)_x \circ (\beta F_1)_x.
    \]
    This is precisely the claim we needed to show, proving the equality. 
\end{remark}

\begin{terminology}\label{A.1.15}
    Let $F,G : \C\to\D$. If there are isomorphisms $F(c)\cong G(c)$ for each $c\in\C$, we say the isomorphism is \emph{natural} in $c$ if they form the components of a natural isomorphism $F(-)\cong G(-)$.
\end{terminology}

\begin{terminology}\label{A.1.16}
    A functor $F$ is called \emph{representable} if it is naturally isomorphic to a hom-functor $\C(c,-)$ or $\C(-,c)=\C\opp(c,-)$.
\end{terminology}
\subsection{Yoneda and universals}\label{secA.2}
A cornerstone of basic category theory is the Yoneda lemma and universal properties. We prove the Yoneda lemma, review the informal notion of a universal property, and the formal definition via limits. Note that we use the notation $\D^\C(-,-)$ for the hom-sets of the functor category. Although this is a standard notation, some authors prefer $\Nat(-,-)$, and this author has found that if you are used to the latter notation, the former may be less intuitive.

\begin{lemma}[Yoneda]\label{A.2.1}
    Let $F : \C \to\Set$ be a functor. Then for all $c\in\C$, there is an isomorphism
    \[
    \Set^\C(\C(c,-),F)\cong F(c)
    \]
    natural in $c$. Dually, let $G : \C\opp\to\Set$ be a functor. Then for all $c\in\C$, there is an isomorphism
    \[
    \Set^{\C\opp}(\C(-,c),G)\cong G(c),
    \]
\end{lemma}
\begin{proof}
    Let $f : c\to d$ in $\C$ and let $\alpha : \C(c,-)\Rightarrow F$. Consider the diagram
    \[\begin{tikzcd}
	{\C(c,c)} & {F(c)} & {1_c} & {\alpha_c(1_c)} \\
	{\C(c,d)} & {F(d)} & f & {\alpha_d(f)=(F(f)\alpha_c)(1_c)}
	\arrow["{\alpha_c}", from=1-1, to=1-2]
	\arrow["{f_*}"', from=1-1, to=2-1]
	\arrow["{F(f)}", from=1-2, to=2-2]
	\arrow[maps to, from=1-3, to=1-4]
	\arrow[maps to, from=1-3, to=2-3]
	\arrow[maps to, from=1-4, to=2-4]
	\arrow["{\alpha_d}"', from=2-1, to=2-2]
	\arrow[maps to, from=2-3, to=2-4]
    \end{tikzcd}\]
    The diagram on the left commutes by naturality of $\alpha$. The diagram on the right follows the image of $1_c$ around the diagram, showing that $\alpha_d(f)$ is uniquely determined by $\alpha_c$ and $F(f)$. Thus the map $\alpha\mapsto \alpha_c(1_c)$ is indeed a bijection. To show the map is natural, consider the functor $\C\to\Set$ given by $c\mapsto \Set^\C(\C(c,-),F)$. The action on morphisms is given by postcomposition: let $f : c\to d$ in $\C$, and let $g : x\to y$ in $\C$. If $\alpha : \C(c,-)\Rightarrow F$, then the square on the left
    \[\begin{tikzcd}
	{\C(d,x)} & {\C(c,x)} & {F(x)} \\
	{\C(d,y)} & {\C(c,y)} & {F(y)}
	\arrow["{f^*}", from=1-1, to=1-2]
	\arrow["{g_*}"', from=1-1, to=2-1]
	\arrow["{\alpha_x}", from=1-2, to=1-3]
	\arrow["{g_*}", from=1-2, to=2-2]
	\arrow["{F(g)}", from=1-3, to=2-3]
	\arrow["{f^*}"', from=2-1, to=2-2]
	\arrow["{\alpha_y}"', from=2-2, to=2-3]
\end{tikzcd}\]
    commutes. We want to show the square on the right also commutes, which will show that $f^* : \C(d,-)\Rightarrow \C(c,-)$, and thus we can define our functor's action on morphisms as $\alpha\mapsto\alpha\circ f^*$. Indeed, for any $h : c\to x$, we have
    \[
    (g_*f^*)(h)=ghf=(f^*g_*)(h).
    \]
    Thus we have our functor. The naturality of the isomorphism thus amounts to showing the square on the left
    \[\begin{tikzcd}
	{\Set^\C(\C(c,-),F)} & {F(c)} & \alpha & {\alpha_c(1_c)} \\
	{\Set^\C(\C(d,-),F)} & {F(d)} & {\alpha f^*} & {(\alpha_d f^*)(1_d)=(F(f)\alpha_c)(1_c)}
	\arrow["{\mathrm{eval}_{1_c}}", from=1-1, to=1-2]
	\arrow["{f^*}"', from=1-1, to=2-1]
	\arrow["{F(f)}", from=1-2, to=2-2]
	\arrow[maps to, from=1-3, to=1-4]
	\arrow[maps to, from=1-3, to=2-3]
	\arrow[maps to, from=1-4, to=2-4]
	\arrow["{\mathrm{eval}_{1_d}}"', from=2-1, to=2-2]
	\arrow[maps to, from=2-3, to=2-4]
\end{tikzcd}\]
    commutes. Following some $\alpha : \C(c,-)\Rightarrow F$ around the square gives the square on the right, so it suffices to show the postulated equality holds. Indeed, by the earlier part of the proof, $(F(f)\alpha_c)(1_c)=\alpha_d(f)$, and we also know $(\alpha_df^*)(1_d)=\alpha_d(f)$. The second statement follows by duality.
\end{proof}

\begin{lemma}[Yoneda]\label{A.2.2}
    For any (locally small) category $\C$, there is a full and faithful functor $y:\C\to\Set^{\C\opp}$, called the \emph{Yoneda embedding}, given by the object function $c\mapsto\C(-,c)$ and the morphism function $f\mapsto f^*$, where $f^*$ denotes precomposition; in particular all natural transformations between representable functors are given by precomposition by a fixed morphism $f$.
\end{lemma}
\begin{proof}
    \Cref{A.2.1} gives us that 
    \[
    \Set^\C(\C(-,c),\C(-,d))\cong\C(c,d),
    \]
    and since we also saw in \cref{A.2.1} that embeddings of this form are functorial, we have our full and faithful embedding. We also saw that the morphism function was $f\mapsto f_*$, and the dual of this is $f\mapsto f^*$, giving us the second part of the statement.
\end{proof}

\begin{definition}\label{A.2.3}
    Let $\C$ be a category. An object $0\in\C$ is \emph{initial} if for all $c\in\C$, there is a \emph{unique} arrow $0\to c$. An object $*\in\C$ is \emph{terminal} if for all $c\in\C$, there is a \emph{unique} arrow $c\to*$.
\end{definition}

\begin{remark}\label{A.2.4}
    Initial and terminal objects need not exist, but if they exist they are unique up to isomorphism. Indeed, let $i,j$ be initial in $\C$. There is a unique arrow $i\to j$ and a unique arrow $j\to i$. Composing these gives arrows $i\to i$ and $j\to j$. Since $i$ and $j$ are initial, there must only be one arrow $i\to i$ and $j\to j$, namely the identity. Thus, the unique arrows are inverses.
\end{remark}

\begin{terminology}\label{A.2.5}
    An object $x$ in a category $\C$ is said to satisfy a \emph{universal property} if it is either initial or terminal \emph{with respect to some property}. This involves constructing a new category of objects satisfying this property, and viewing the object as initial or terminal in it. For example, fix some $c\in\C$, and suppose the property is that an object $x\in\C$ has a specified arrow $f : x\to c$. Construct the category of pairs $(x,f)$ where $f : x\to c$, and define morphisms $\varphi : (f,x)\to (g,y)$ as morphisms $\varphi : x\to y$ in $\C$ such that the diagram
    \[\begin{tikzcd}
	x & y \\
	c
	\arrow["\varphi", from=1-1, to=1-2]
	\arrow["f"', from=1-1, to=2-1]
	\arrow["g", from=1-2, to=2-1]
\end{tikzcd}\]
    commutes. An initial object in this category is thus a pair $(h,u)$ such that for any other $(f,x)$ in the category, there is a unique morphism $(h,u)\rightarrow (f,x)$. In particular, there is a unique morphism $u\to x$ making the diagram
    \[\begin{tikzcd}
	u & x \\
	c
	\arrow["\varphi", dashed, from=1-1, to=1-2]
	\arrow["h"', from=1-1, to=2-1]
	\arrow["f", from=1-2, to=2-1]
\end{tikzcd}\]
    commute. We denote unique morphisms with dashed arrows. The pair $(h,u)$ is then said to satisfy this universal property. Univerals need not exist, but if they do, they are unique up to isomorphism.
\end{terminology}

\begin{definition}\label{A.2.6}
    Let $F : \J\to\C$ be a functor. A \emph{cone} for $F$ is an object $c\in C$ and a collection of morphisms $\{\mu_j\}_{j\in\J}$ for each object of $\J$, such that for any morphism $f : j_1\to j_2$ in $\J$, the diagram
    \[\begin{tikzcd}
	& c \\
	{F(j_1)} && {F(j_2)}
	\arrow["{\mu_1}"', from=1-2, to=2-1]
	\arrow["{\mu_2}", from=1-2, to=2-3]
	\arrow["{F(f)}"', from=2-1, to=2-3]
\end{tikzcd}\]
    commutes. If $(c,\mu)$ and $(d,\nu)$ are cones for $F$, a \emph{morphism} of cones $\varphi : (d,\nu)\to(c,\mu)$ is an arrow $\varphi : d\to c$ such that for all $f : j_1\to j_2$ in $\J$, the diagram
    \[\begin{tikzcd}
	& d \\
	& c \\
	{F(j_1)} && {F(j_2)}
	\arrow["{\nu_1}"', curve={height=12pt}, from=1-2, to=3-1]
	\arrow["{\nu_2}", curve={height=-12pt}, from=1-2, to=3-3]
	\arrow["{\mu_1}"', from=2-2, to=3-1]
	\arrow["{\mu_2}", from=2-2, to=3-3]
	\arrow["{F(f)}"', from=3-1, to=3-3]
    \arrow["\varphi", from=1-2, to=2-2]
\end{tikzcd}\]
    commutes.

    Dually, a \emph{co-cone} for $F$ is an object $c\in\C$ and a collection of morphisms $\{\mu_j\}_{j\in\J}$ for each object of $\J$, such that for any morphism $f : j_1\to j_2$ in $\J$, the dual diagram
    \[\begin{tikzcd}
	{F(j_1)} && {F(j_2)} \\
	& c
	\arrow["{F(f)}", from=1-1, to=1-3]
	\arrow["{\mu_1}"', from=1-1, to=2-2]
	\arrow["{\mu_2}", from=1-3, to=2-2]
\end{tikzcd}\]
    commutes. If $(c,\mu)$ and $(d,\nu)$ are co-cones for $F$, a \emph{morphism} of co-cones $\varphi : (d,\nu)\to(c,\mu)$ is an arrow $\varphi : d\to c$ such that for all $f : j_1\to j_2$ in $\J$, the diagram
    \[\begin{tikzcd}
	{F(j_1)} && {F(j_2)} \\
	& d \\
	& c
	\arrow["{F(f)}", from=1-1, to=1-3]
	\arrow["{\nu_1}"', from=1-1, to=2-2]
	\arrow["{\mu_1}"', curve={height=12pt}, from=1-1, to=3-2]
	\arrow["{\nu_2}", from=1-3, to=2-2]
	\arrow["{\mu_2}", curve={height=-12pt}, from=1-3, to=3-2]
	\arrow["\varphi", from=2-2, to=3-2]
\end{tikzcd}\]
    commutes.
\end{definition}

\begin{definition}\label{A.2.7}
    Let $F:\J\to\C$ be a functor. A \emph{limit} for $F$ is a universal cone $(\lim F,\mu)$ such that for any other cone $(c,\nu)$ for $F$, there is a unique dashed arrow
    \[\begin{tikzcd}
	& c \\
	& {\lim F} \\
	{F(j_1)} && {F(j_2)}
	\arrow[dashed, from=1-2, to=2-2]
	\arrow["{\nu_1}"', curve={height=12pt}, from=1-2, to=3-1]
	\arrow["{\nu_2}", curve={height=-12pt}, from=1-2, to=3-3]
	\arrow["{\mu_1}"', from=2-2, to=3-1]
	\arrow["{\mu_2}"', from=2-2, to=3-3]
	\arrow["{F(f)}"', from=3-1, to=3-3]
\end{tikzcd}\]
    rendering the diagram commutative for all $f : j_1\to j_2$ in $\J$.

    Dually, a \emph{colimit} for $F$ is a universal co-cone $(\colim F,\mu)$ such that for any other co-cone $(c,\nu)$ for $F$, there is a unique dashed arrow
    \[\begin{tikzcd}
	{F(j_1)} && {F(j_2)} \\
	& {\colim F} \\
	& c
	\arrow["{F(f)}", from=1-1, to=1-3]
	\arrow["{\mu_1}"', from=1-1, to=2-2]
	\arrow["{\nu_1}"', curve={height=12pt}, from=1-1, to=3-2]
	\arrow["{\mu_2}", from=1-3, to=2-2]
	\arrow["{\nu_2}", curve={height=-12pt}, from=1-3, to=3-2]
	\arrow[dashed, from=2-2, to=3-2]
\end{tikzcd}\]
    rendering the diagram commutative for all $f : j_1\to j_2$ in $\J$.
\end{definition}

\begin{remark}\label{A.2.8}
    In the above definitions, $\J$ is usually a small category, called an \emph{indexing} category. The functor $F$ is often called a \emph{diagram}. The intuition is that the functor $F$ indexes a diagram in $\C$. Many authors define universal properties as (co)limits of functors, or as an equivalent notion. Limits and colimits, like any universals, need not exist, but if they exist they are unique up to isomorphism.
\end{remark}

\begin{terminology}\label{A.2.9}
    Let $\J$ be a category. If a category $\C$ has all limits of functors $\J\to\C$, we call it $\J$-complete. If $\C$ has all colimits of functors $\J\to\C$, we call it $\J$-cocomplete. If $\C$ has all limits of all small diagrams (functors with small domain), we call $\C$ \emph{complete}. If $\C$ has all colimits of all small diagrams, we call $\C$ \emph{cocomplete}. We also say $\C$ has all small (co)limits ( respectively finite (co)limits, $\J$-indexed (co)limits) when the respective notion exists.
\end{terminology}

\begin{proposition}\label{A.2.10}
    Let $\C$ have limits over $\J$. Then $\lim : \C^\J\to\C$ is a functor. Dually, if $\C$ has colimits over $\J$, then $\colim : \C^\J\to\C$ is a functor.
\end{proposition}
\begin{proof}
    The object function is given. Let $F,G : \J\to\C$, and let $\alpha : F\Rightarrow G$. Componentwise postcomposition of the limiting cone for $F$ with $\alpha$ gives a cone for $G$ by naturality, and by universality of limits, there is a unique induced arrow $\lim F\to\lim G$ making the relevant diagram commute. Define the induced arrow as $\lim\alpha$.
\end{proof}
\subsection{Products and coproducts}\label{secA.3}
Products in arbitrary categories, and the dual notion of coproducts, are constructions of fundamental importance to the field of category theory. For example, any limit can be constructed using only products and equalizers, and any colimit can be constructed using only coproducts and coequalizers \cite[V.2]{ml}. Any category with finite products may be regarded as a symmetric monoidal category, which is a particularly important class of monoidal categories. This section will introduce (co)products from the perspective of universal properties and (co)limits, will give some simple examples, and will develop the product and coproduct bifunctors.

\begin{definition}\label{A.3.1}
    Let $\C$ be a category, and let $x,y\in\C$. The \emph{product} of $x$ and $y$, denoted $x\times y$, is an object of $\C$, together with two morphisms $\pi_x : x\times y \to x$ and $\pi_y : x\times y \to y$, which are universal in the sense that for any diagram of solid arrows
    \[\begin{tikzcd}[row sep=3.5em, column sep=3.5em]
	& z \\
	x & {x\times y} & y
	\arrow["f"', from=1-2, to=2-1]
	\arrow["{(f,g)}"', dashed, from=1-2, to=2-2]
	\arrow["g", from=1-2, to=2-3]
	\arrow["{\pi_x}", from=2-2, to=2-1]
	\arrow["{\pi_y}"', from=2-2, to=2-3]
    \end{tikzcd}\]
    there is a unique dashed arrow as indicated, rendering the diagram commutative.
\end{definition}

\begin{definition}\label{A.3.2}
    Let $\C$ be a category, and let $x,y\in\C$. The \emph{coproduct} of $x$ and $y$, denoted $x\amalg y$, is the dual notion of a product, meaning it is an object of $\C$ together with two morphisms $i_x : x \to x\amalg y$ and $i_y : y\to x\amalg y$, which are universal in the sense that for any diagram of solid arrows
    \[\begin{tikzcd}[row sep=3.5em, column sep=3.5em]
	x & {x\amalg y} & y \\
	& z
	\arrow["{i_x}", from=1-1, to=1-2]
	\arrow["f"', from=1-1, to=2-2]
	\arrow["{f+g}", dashed, from=1-2, to=2-2]
	\arrow["{i_y}"', from=1-3, to=1-2]
	\arrow["g", from=1-3, to=2-2]
    \end{tikzcd}\]
    there is a unique dashed arrow as indicated, rendering the diagram commutative.
\end{definition}

\begin{notation}\label{A.3.3}
    As indicated in \cref{A.3.1}, we typically denote the unique arrow induced by $f$ and $g$ as $(f,g)$. As indicated indicated in \cref{A.3.2}, we typically denote the unique arrow induced by $f$ and $g$ as $f+g$. The arrows $\pi_x$ and $\pi_y$ are called the \emph{natural projections}, and the arrows $i_x,i_y$ are called the \emph{natural injections}. We will often drop the ``natural'' prefix and refer to the $\pi$s as projections and the $i$s as injections, but for some reason we drop the prefix more often for injections than projections.
\end{notation}

\begin{terminology}\label{A.3.4}
    Recall that a \emph{discrete} category is a category $\J$ which has only identity arrows. We oftentimes will regard as set $J$ as a category $\J$ by taking $\J$ to be the discrete category whose objects are the elements of $J$. Since $\J$ offers no additional structure beyond the objects, this is a fair definition. We oftentimes abuse language and refer to a set as if it were a category, with the understanding that we are really referring to the discrete category that represents the set.
\end{terminology}

\begin{remark}\label{A.3.5}
    Let $\J$ be the set $\{0,1\}$ of 2 elements. A functor $F : \J \to \C$ is simply a choice of two elements $F(0)=x$ and $F(1)=y$, since there are no nontrivial morphisms to consider. A cone for $F$ is thus an object $c\in\C$ and arrows $\mu_0,\mu_1$ as described by the diagram
    \[\begin{tikzcd}
	& c \\
	{F(0)} && {F(1)}
	\arrow["{\mu_0}", from=1-2, to=2-1]
	\arrow["{\mu_1}"', from=1-2, to=2-3]
    \end{tikzcd}\]
    A \emph{limit} of $F$ is thus an object $\lim F$ of $\C$, together with two morphisms $\mu_0 : \lim F \to F(0)$ and $\mu_1 : \lim F \to F(1)$, such that for any cone $(c,\nu)$ for $F$, there is a unique arrow $c\to \lim F$ as indicated in the diagram
    \[\begin{tikzcd}[row sep=3.5em, column sep=3.5em]
	& c \\
	F(0) & {\lim F} & F(1)
	\arrow["\nu_0"', from=1-2, to=2-1]
	\arrow[dashed, from=1-2, to=2-2]
	\arrow["\nu_1", from=1-2, to=2-3]
	\arrow["{\mu_0}", from=2-2, to=2-1]
	\arrow["{\mu_1}"', from=2-2, to=2-3]
    \end{tikzcd}\]
    rendering it commutative. Note that this is precisely the universal property of \cref{A.3.1}, so the product of $x$ and $y$ is simply the limit of the functor $F : \J\to\C$ such that $F(0)=x$ and $F(1)=y$. By duality, the coproduct $x\amalg y$ is the colimit $\colim F$.
\end{remark}

\begin{definition}\label{A.3.6}
    Let $\C,\D$ be categories. The product category $\C\times\D$ is defined as the category where
    \begin{itemize}
        \item objects are ordered pairs $(c,d)$ with $c\in\C$ and $d\in\D$;
        \item morphisms $(c,d)\to (c',d')$ are ordered pairs $(f,g)$ such that $f : c\to c'$ in $\C$ and $g : d\to d'$ in $\D$.
    \end{itemize}
\end{definition}

\begin{construction}\label{A.3.7}
    A functor whose domain is a product category is called a \emph{bifunctor}. If a category $\C$ has all binary products --- meaning for any $c,d\in\C$, the product $c\times d$ exists --- then product assembles into a bifunctor $-\times - : \C\times\C\to\C$. The action on objects is clear: a pair $(c,d)\in\C\times\C$ is mapped to its product $c\times d$ in $\C$. The action on morphisms is a little more elusive, so consider a morphism $(f,g) : (c,d) \to (c',d')$ in $\C\times\C$; this is equivalently two morphisms
    \[\begin{tikzcd}[row sep=0em]
	c & {c'} \\
	d & {d'}
	\arrow["f", from=1-1, to=1-2]
	\arrow["g"', from=2-1, to=2-2]
    \end{tikzcd}\]
    in $\C$. We can construct the diagram of solid arrows
    \[\begin{tikzcd}
	c && {c'} \\
	{c\times d} && {c'\times d'} \\
	d && {d'}
	\arrow["f", from=1-1, to=1-3]
	\arrow["{\pi_c}", from=2-1, to=1-1]
	\arrow["{f\times g}", dashed, from=2-1, to=2-3]
	\arrow["{\pi_d}"', from=2-1, to=3-1]
	\arrow["{\pi_{c'}}"', from=2-3, to=1-3]
	\arrow["{\pi_{d'}}", from=2-3, to=3-3]
	\arrow["g"', from=3-1, to=3-3]
    \end{tikzcd}\]
    which commtues manifestly. Note that the composites $f\circ \pi_c$ and $g\circ \pi_d$ are arrows from $c\times d$ to $c'$ and $d'$, respectively, so by universality of product there is a unique morphism $c\times d \to c'\times d'$ as indicated by the dashed arrow, rendering the diagram commtutative. We define $f\times g$ as that unique arrow. Functoriality is easily shown by universality. We similarly construct the coproduct bifunctor $-\amalg - : \C\times\C \to\C$.
\end{construction}

\begin{remark}\label{A.3.8}
    Using the definition of a product in \cref{A.3.5}, we can more easily define arbitrary products. Given a set $\J$, a $\J$-indexed product in $\C$ is a limit of a functor $F : \J \to \C$. This functor selects a collection $\{x_j\}_{j\in\J}$ of objects in $\C$, and taking the limit corresponds to taking the product over this collection. This allows us to even define limits where the indexing object is not a set, but a class, by viewing the class as a large discrete category. We call products over large categories \emph{large products}, products over small categories (sets) are \emph{small products}, and products over finite categories (finite sets) are \emph{finite products}. The same holds for coproducts.
\end{remark}

\begin{example}\label{A.3.9}
    Products in $\Set$ are our usual cartesian products of sets. Coproducts in $\Set$ are the disjoint union: we take two sets $X$ and $Y$, and forcibly make them disjoint by replacing $X$ with $X\times\{0\}$ and $Y\times\{1\}$, and then take the union of these new sets. Note that the choices of 0 and 1 here are completely arbitrary, this is simply one way to force the sets to be disjoint, and any other method will work just as well. The only think that matters is that all options are isomorphic to each other in $\Set$, and indeed universals are unique up to isomorphism. Some other examples are
    \begin{itemize}
        \item The product in $\Grp$ is the direct product of groups, and the coproduct is the free product;
        \item The product in $\Ab$ is the direct product of groups, but the (finite) coproduct is \emph{also} the direct product of groups;
        \item Just line in $\Ab$, the (finite) product and coproduct in any category $\Mod_R$ of $R$-modules coencide, and indeed this happens in any \emph{abelian category} \cite[VIII.2]{ml}, \cite[IX.1.4]{alg0}. In abelian categories, we call the (finite) product/coproduct a \emph{biproduct}, and in $\Ab$ and $\Mod_R$ we call it the \emph{direct sum}. We typically denote it with the symbol $\oplus$;
        \item The product in the poset category $\mathbb{Z}$ of a pair $(n,m)$ is the minimum number in the pair, and the coproduct is the maximum;
        \item The product in $\Top$ is the product of topological spaces, and the coproduct is the disjoint union, with the usual topologies.
    \end{itemize}
\end{example}

\begin{construction}\label{A.3.10}
    The product of two categories as given in \cref{A.3.6} indeed satisfies the universal property of \cref{A.3.1}, so we obtain by way of \cref{A.3.7} a product bifunctor on the category $\Cat$ of small categories. The reason we restrict here to the category of \emph{small} categories has nothing to do with products, but with the foundational issues that arise when trying to define a category of all categories. However, we can still define the product of two functors, even if they don't live between small categories. They are constructed in the exact way you would expect.

    Let $F : \C\to\D$ and $G : \C'\to\D'$ be functors. Then the product functor $F\times G : C\times C' \to \D\times\D'$ is defined on objects as
    \[
    (F\times G)(c,c')=(F(c),G(c'))
    \]
    and on morphisms as
    \[
    (F\times G)(f,g)=(F(f),G(g)).
    \]
\end{construction}
\subsection{Adjunctions}\label{secA.4}
Mac Lane, one of the originators of category theory, is famously quoted in \cite{ml} as saying ``adjoint functors arise everywhere.'' Indeed, the notion of adjunction is ever present in category theory and beyond, and we require our reader to be familiar with the concept. In this section, we give the equivalent definitions via hom-sets and unit-counits, and we introduce some basic results about adjoint functors, such as preservation of limits and pointwise construction from comma categories or universal arrows. We do not prove these results; the reader interested in a more rigorous approach should consult \cite{ctx}, \cite{lein}, or \cite{ml}. Most of our presentation follows that of \cite{ml}.

\begin{definition}\label{A.4.1}
    Let $F : \C\to\D$ and $U:\D\to\C$ be functors. An \emph{adjunction} between $F$ and $U$ is an isomorphism
    \[
    \varphi_{c,d} : \D(F(c),d)\cong \C(c,U(d))
    \]
    for each $c\in\C$ and $d\in\D$, natural in both variables.
\end{definition}

\begin{notation}\label{A.4.2}
    There are a lot of important notations concerning adjoint functors, and different authors will use various different notations. If $\varphi$ is an adjunction
    \[
    \varphi_{c,d} : \D(F(c),d)\cong \C(c,U(d)),
    \]
    then we say $F$ is \emph{left adjoint} to $U$, and $U$ is \emph{right adjoint} to $F$, since $F$ shows up on the left slot of the hom-set, and $U$ shows up on the right. We write $F\dashv U : \C\to\D$ to denote that $F$ is left adjoint to $U$, and that $F:\C\to\D$ and $U:\D\to\C$. When the categories can be inferred from context, or when writing them is unnecessary, we drop the categories and write only $F\dashv U$. We also write $(F,U,\varphi):\C\to\D$ to denote that $\varphi$ is an adjunction $F\dashv U$, where $F : \C\to\D$ and $U:\D\to\C$. Once again, if the categories can be inferred from context, or if writing them is unnecessary, we drop the categories and instead write $(F,U,\varphi)$. If $F\dashv U$ is an adjunction, we say $F$ is a \emph{left adjoint functor}, and $U$ is a \emph{right adjoint functor}, since $F$ is left adjoint to $U$, and vice versa. If we say ``$F : \C \rightleftarrows \D : U$ is an adjunction'', we always take $F$ to be left adjoint to $U$. We may also write
    \[\begin{tikzcd}
	\C & \D
	\arrow[""{name=0, anchor=center, inner sep=0}, "F", shift left=2, from=1-1, to=1-2]
	\arrow[""{name=1, anchor=center, inner sep=0}, "U", shift left=2, from=1-2, to=1-1]
	\arrow["\dashv"{anchor=center, rotate=-90}, draw=none, from=0, to=1]
\end{tikzcd}\]
    to denote that $F\dashv U$. When we say the isomorphism $\varphi$ is \emph{natural in both variables}, we mean that $\varphi$ is a natural isomorphism
    \[
    \varphi:\D(F(-),-)\cong \C(-,U(-)).
    \]
\end{notation}

\begin{proposition}\label{A.4.3}
    The definition of an adjunction as given in \cref{A.4.1} is equivalent to the following definition: two functors $F : \C\to\D$ and $U:\D\to\C$ are \emph{adjoint} if there is a natural transformation $\eta : 1_\C\Rightarrow UF$ called the \emph{unit of the adjunction}, and a natural transformation $\varepsilon : FU\Rightarrow 1_\D$ called the \emph{counit of the adjunction}, such that the the diagrams
    \[\begin{tikzcd}
	U & UFU & FUF & F \\
	& U & F
	\arrow["{\eta U}", from=1-1, to=1-2]
	\arrow[equals, from=1-1, to=2-2]
	\arrow["{U\varepsilon }", from=1-2, to=2-2]
	\arrow["{\varepsilon F}"', from=1-3, to=2-3]
	\arrow["{F\eta }"', from=1-4, to=1-3]
	\arrow[equals, from=1-4, to=2-3]
    \end{tikzcd}\]
    commute. These are known as the \emph{triangle identities}, and they are equivalently written as
    \[
    1_F = \varepsilon F\circ F\eta,\qquad 1_U = U\varepsilon\circ\eta U.
    \]
    For any $c\in\C$ and $d\in\D$, the components $\eta_c$ and $\varepsilon_d$ of the unit and counit are given as
    \[
    \eta_c = \varphi(1_{F(c)}),\qquad \varepsilon(d)=\varphi^{-1}(1_{U(d)}),
    \]
    where $\varphi$ is the adjunction as given in \cref{A.4.1}. This provides a way to construct the unit and counit given an adjunction $(F,U,\varphi):\C\to\D$, or a way to construct the adjunction $\varphi$ given a unit $\eta$ and a counit $\varepsilon$. It is also shown in \cite{ml} that each $(\eta_c,F(c))$ is a universal arrow from $c$ to $U$, and each $(\varepsilon_d,U(d))$ is a universal arrow from $F$ to $d$. We do not require the reader to know what this means, but it is an interesting aside.
\end{proposition}

\begin{notation}\label{A.4.4}
    We write $(F,U,\eta,\varepsilon) : \C\to\D$ to denote that there is an adjunction $F\dashv U$ with $\eta$ the unit and $\varepsilon$ the counit of the adjunction. When the categories can be inferred from context, or when it is unnecessary to write them, we instead write $(F,U,\eta,\varepsilon)$.
\end{notation}

\begin{proposition}\label{A.4.5}
    If $F$ is a left adjoint functor, then it preserves colimits. If $U$ is a right adjoint functor, then it preserves limits. More precisely, if $F\dashv U : \C\to\D$ is an adjunction, and $G:\J\to\C$ and $H:\J\to\D$ are diagrams (with $\J$ not necessarily small), then
    \[
    F(\colim G)\cong \colim (F\circ G),\qquad U(\lim H)\cong \lim(U\circ H).
    \]
\end{proposition}

\begin{remark}\label{A.4.6}
    Let $U:\D\to\C$ be a functor such that for all $c\in\C$, the comma category $(c\downarrow U)$ has an initial object $u : c\to U(r)$. This means for any $f : c\to U(x)$, there is a unique morphism $r\to x$ in $\D$ such that the diagram
    \[\begin{tikzcd}
	c & {U(r)} & r \\
	& {U(x)} & x
	\arrow["u", from=1-1, to=1-2]
	\arrow["f"', from=1-1, to=2-2]
	\arrow["{U(g)}", dashed, from=1-2, to=2-2]
	\arrow["g", dashed, from=1-3, to=2-3]
    \end{tikzcd}\]
    commutes. An interesting aside is that this means $(r,u)$ is a universal arrow from $c$ to $U$. We construct a functor $F:\C\to\D$ as follows:
    \begin{enumerate}
    \item Let $c\in\C$, and let $u:c\to U(r)$ be the initial object in $(c\downarrow U)$. Define the object function of $F$ by mapping $c\mapsto r$. One might argue this is ill-defined, since $U$ need not be injective, but objects in the comma category are actually pairs $(u,r)$, not $(u,U(r))$, so the datum of an object in the comma category includes the specific preimage under $U$ we are considering. One then needs to only make a specific choice of which initial object to choose, which follows by invoking the Axiom of Choice.
        \item Let $c_1,c_2\in C$, and let $f : c_1\to c_2$. We define the morphism function of $F$ as follows. Let $u_1 : c_1\to U(r_1)$ be the initial object in $(c_1\downarrow U)$, and likewise $u_2:c_2\to U(r_2)$ the initial object in $(c_2\downarrow U)$. The composite $u_2\circ f : c_1\to U(r_2)$ induces a diagram of solid arrows
        \[\begin{tikzcd}
	{c_1} & {U(r_1)} & {r_1} \\
	{c_2} & {U(r_2)} & {r_2}
	\arrow["{u_1}", from=1-1, to=1-2]
	\arrow["f"', from=1-1, to=2-1]
	\arrow["{U(g)}", dashed, from=1-2, to=2-2]
	\arrow["g", dashed, from=1-3, to=2-3]
	\arrow["{u_2}"', from=2-1, to=2-2]
    \end{tikzcd}\]
    Since $u_2\circ f : c_1\to U(r_2)$ is an object in $(c_1\downarrow U)$, and since $u_1 : c_1\to U(r_1)$ is initial in this category, there exists a dashed arrow $g : r_1\to r_2$ rendering the diagram commutative. Define the morphism function of $F$ as $f\mapsto g$. This is indeed well-defined by uniqueness of $g$, and functoriality follows easily from universality.
    \end{enumerate}
\end{remark}

\begin{proposition}\label{A.4.7}
    The functor $F$ constructed in \cref{A.4.6} is left adjoint to $U$. Dually, if we are given $F:\C\to\D$ and each $(F\downarrow d)$ has a terminal object, one may assemble a functor $U:\D\to\C$ which is right adjoint to $F$ in the same manner.
\end{proposition}